% ── Rakouský model: Kammer-system a Allgemeinverbindlicherklärung ────────────

Rakouský model je~v~evropském srovnání založen na~silném odvětvovém vyjednávání
a~vysoké organizovanosti zaměstnanců i~zaměstnavatelů díky povinnému členství v~hospodářské
komoře (\textit{Wirtschaftskammer Österreich, WKÖ}) a~komoře pracovníků (\textit{Arbeiterkammer, AK}).
Členství v~odborových organizacích (největší je sdružení \textit{Österreichischer Gewerkschaftsbund, ÖGB}) není povinné.~\cite{Eurofound_AT_Actors2024}\par
V~praxi to znamená, že pokrytí \acs{KS} zůstává dlouhodobě velmi vysoké i~při nižší odborové
organizovanosti, než je běžné v~severských zemích. Podniková úroveň zde typicky
plní často implementační roli vůči odvětvovým standardům.\par

V~datech \acs{OECD}~\acs{ICTWSS} se~Rakousko dlouhodobě drží u~hodnot pokrytí
blízkých \SI{98}{\percent} (vyšší v~soukromé sféře), zatímco
odborová organizovanost je výrazně nižší (přibližně \SI{26}{\percent}).
Tento případ ukazuje, že vysokého pokrytí \acs{KS} lze dosáhnout s nižší odborovou organizovaností.~\cite{Eurofound_AT_CB2024}~\cite{OECD_UnionEmployerCoverage_2025}\par

Vysoké pokrytí zároveň snižuje prostor pro~vnitrostátní mzdový dumping mezi
podniky téhož odvětví, protože minimální mzdy jsou stanovovány odvětvovými smlouvami.
To je rozdíl vůči \acloc{ČR}, kde je odvětvová vrstva velmi slabá a~celkové pokrytí \acs{KS} nízké.
Minimální mzda stanovená zákonem pro všechny sektory hospodářství neodráží odvětvové rozdíly, které se prohlubují.\par


Rakouský model představuje pro~\acacc{ČR} nejbližší inspirační rámec:
geograficky sousedící ekonomiku s~podobnou strukturou, ale s~výrazně
silnějším odvětvovým vyjednáváním. Pro českou debatu to znamená, že těžiště
reforem neleží pouze v~posilování odborové a zaměstnavatelské organizovanosti, ale v~praktické
aktivaci existujícího nástroje \S~7 zákona č.~2/1991~Sb., v~datové podpoře
reprezentativity a~v~systémové podpoře sektorového vyjednávání.\par
